\documentclass[11pt,a4paper]{report}

\usepackage[margin=1in]{geometry}
\usepackage{listings}
\usepackage{xcolor}
\usepackage{booktabs}
\usepackage{longtable}
\usepackage{hyperref}
\usepackage{caption}
\usepackage{amsmath}
\usepackage{enumitem}
\usepackage{fancyhdr}

\hypersetup{colorlinks=true, linkcolor=blue, urlcolor=blue, citecolor=blue}

\definecolor{codebg}{rgb}{0.96,0.96,0.96}
\definecolor{codekw}{rgb}{0.0,0.0,0.6}
\definecolor{codecomment}{rgb}{0.3,0.5,0.3}
\definecolor{codestring}{rgb}{0.6,0.0,0.0}

\lstdefinestyle{cppstyle}{
  language=C++,
  backgroundcolor=\color{codebg},
  basicstyle=\ttfamily\small,
  commentstyle=\color{codecomment}\itshape,
  keywordstyle=\color{codekw}\bfseries,
  stringstyle=\color{codestring},
  numbers=left,
  numberstyle=\tiny\color{gray},
  numbersep=8pt,
  frame=single,
  rulecolor=\color{gray!40},
  breaklines=true,
  breakatwhitespace=false,
  showstringspaces=false,
  tabsize=4,
  columns=fullflexible,
  keepspaces=true
}
\lstset{style=cppstyle}

\pagestyle{fancy}
\fancyhf{}
\fancyhead[L]{APARA Simulator (engine\_isp) --- Bug Report}
\fancyhead[R]{\thepage}
\renewcommand{\headrulewidth}{0.4pt}

\title{\textbf{APARA Simulator (\texttt{engine\_isp}) Bug Report \\ Solved and Unsolved Defects, With Source}}
\author{APARA C Compiler Project --- Verification \& Audit}
\date{\today}

\begin{document}

\maketitle
\tableofcontents

% ============================================================================
\chapter{Introduction and Methodology}
% ============================================================================

This report documents every confirmed defect found in the APARA hardware simulator's own
source code (\texttt{engine\_isp}'s C++ assembler/simulator, \emph{not} the Python C
compiler being developed alongside it) across the full history of this project. Each entry
was confirmed by checking that the compiler emitted the textually-correct mcode instruction
for the case in question, and that the simulator's own assembly or execution of that
\emph{correct} instruction is what produced the wrong result --- ruling out a compiler bug
before ever blaming the simulator.

\section{Scope}
Six confirmed defects are documented: three have been found, root-caused, fixed, and
verified as currently present in the authoritative simulator binary; three remain
confirmed but unfixed, and are flagged here for the course staff/hardware team, since they
sit outside what a C-compiler project can fix on its own. Three further findings are
\emph{not} source-code bugs but are directly relevant environment/tooling issues
encountered while building this verification: two disagreeing simulator binaries in the
project tree, a divergence between two simulator source checkouts, and a hard
architectural limitation (not a bug) blocking function pointers.

\section{Authoritative source and binary}
All line numbers and code listings in this report are read directly from the source tree
backing the project's authoritative simulator binary, \texttt{mcode\_run}, confirmed by
matching file timestamps and verified against the project's own build records. All
``current status'' claims (fixed/not fixed, present/absent) were re-checked against this
live source immediately before writing this report --- not assumed from memory.

\section{How these were found}
The three solved bugs (Chapter~\ref{ch:solved}) surfaced from a standalone investigation
into nested function calls returning garbage: any C function called from within another
function, where the callee was defined \emph{before} its caller in source order (the
ordinary C pattern), returned a stale or zero value instead of computing correctly. Six
separate compiler-side hypotheses --- bundle padding, jump-target resolution, bundle
shape/ordering, caller-save/restore ordering, callee-internal codegen, and epilogue
scratch-register clobber --- were checked and individually ruled out before the
investigation moved into the simulator's own source and found the real causes there.

The three unsolved bugs (Chapter~\ref{ch:unsolved}) surfaced from a systematic
instruction-coverage sweep with an independent ``no bias'' verification methodology: every
test's expected result is computed by compiling the \emph{same} source natively with
\texttt{gcc} against a plain-C reference implementation, never by this project's own
compiler or by hand-derivation, and checked against the simulator's real execution via its
own built-in \texttt{PostCondition} mechanism.

% ============================================================================
\chapter{Solved Bugs}
\label{ch:solved}
% ============================================================================

These three bugs were root-caused, fixed, and the fixes are confirmed present in the
\emph{current} authoritative simulator source and binary --- not just a one-off
verification build. All three fixes are exercised by tests that are part of the project's
current, continuously-passing regression baseline.

\section{Bug E1: \texttt{\$ld} sub-word loads ignored the real byte offset, always reading the upper half-word}

\textbf{Status}: \textbf{SOLVED} (fixed upstream by the hardware/VM team; fixed binaries
pulled into this project's simulator build; hardware-confirmed).

\subsection*{Symptom}
\texttt{\$ld (\$i32)} always returned bits \texttt{[63:32]} (the upper 32 bits) of the
8-byte DMEM word being addressed, regardless of the byte offset actually encoded in the
load instruction. \texttt{\$st (\$i32)} was correct --- it wrote to the upper or lower 32
bits depending on the real offset --- only the \emph{load} side ignored the offset. Any
value stored at byte-offset 4 of a word (the conventional placement for a value that isn't
the first element in that word) could never be read back correctly by a 32-bit load.

\subsection*{Why this mattered}
This is the load instruction underlying essentially every \texttt{int}, \texttt{short}, and
\texttt{char} scalar or array access compiled from C --- i.e.\ almost everything a real
program does.

\subsection*{Workaround used while the bug was live (compiler-side, no longer needed)}
Every C element, regardless of type, was given its own full 8-byte-aligned DMEM slot
(\texttt{stride = max(elem\_bytes, 8)}), and the compiler always emitted \texttt{(\$i64)}
loads/stores regardless of the real element width --- sidestepping the broken sub-word
addressing path entirely rather than relying on it. The compiler's narrower load/store
support was only re-enabled \emph{after} this fix was confirmed on hardware.

\subsection*{Fix and verification}
Fixed upstream in the VM/hardware build; the corrected binaries were pulled into this
project's \texttt{engine\_isp} checkout. A full 19-program hardware regression
(\texttt{test\_subword\_i8.c}, \texttt{test\_subword\_i16.c}, isolated single-width probes)
both returned \texttt{r1=1} (full pass). \texttt{test\_array}, \texttt{test\_cmov},
\texttt{test\_branch}, and \texttt{test\_pointer} --- all exercising \texttt{\$i32} in other
shapes --- were independently confirmed correct on the same hardware run. This fix
underlies the compiler's current \texttt{i8}/\texttt{i16}/\texttt{i32} sub-word load/store
support and is part of the project's continuously-passing baseline.

\section{Bug E2: \texttt{\$call}'s disassembler sign-extended the wrong bit, breaking every backward function call}

\textbf{Status}: \textbf{SOLVED} --- confirmed present in the current source (verified by
direct inspection immediately before writing this report).

\subsection*{The code, as currently fixed (\texttt{McodeDisassemble.cpp}, function \texttt{DisassembleToCallInstr})}

\begin{lstlisting}[caption={\texttt{McodeDisassemble.cpp}, lines 261--284 (current, fixed state)}]
McodeInstruction* DisassembleToCallInstr (uint32_t pc, uint32_t hex_instr)
{
    uint32_t opcode  = Get_Slice(31, 26, hex_instr);
    uint32_t imm_flag = Get_Slice (25,25, hex_instr);
    uint32_t rs2  = Get_Slice (4,0, hex_instr);
    int32_t relative_jump = (int32_t) Sign_Extend(24, Get_Slice (24, 0, hex_instr));
    // FIX: was Sign_Extend(25, ...) -- see explanation below

    McodeCallInstruction* ci = NULL;
    if(imm_flag)
    {
        int32_t label_address = ((int32_t) pc)  + relative_jump;
        assert(label_address >= 0);

        string lbl  = "l_" + Int64ToString(label_address);
        ci = new McodeCallInstruction (lbl);
        ci->Set_Pc_Relative_Jump_Address (relative_jump);
    }
    else
    {
        ci = new McodeCallInstruction (rs2);
    }
    return(ci);
}
\end{lstlisting}

\subsection*{The bug (before this fix)}
The line read \texttt{Sign\_Extend(25, Get\_Slice(24, 0, hex\_instr))}. \texttt{Sign\_Extend}'s
first argument is a \emph{bit index}, defined as (\texttt{McodeUtils.cpp}, line 371):

\begin{lstlisting}[caption={\texttt{McodeUtils.cpp}, lines 371--382 --- \texttt{Sign\_Extend}'s definition}]
uint32_t Sign_Extend (uint32_t sign_index, uint32_t x)
{
    bool pad_ones = (((1 << sign_index) & x) != 0);
    uint32_t or_mask = ((~0) << (sign_index+1));

    if(sign_index >= 31)
        or_mask = 0;

    uint32_t ret_val = pad_ones ? (x | or_mask) : x;
    return(ret_val);
}
\end{lstlisting}

\texttt{\$call}'s jump field is 25 bits wide --- bits \texttt{24:0} --- so its real sign bit
sits at index \textbf{24}, not 25. \texttt{Get\_Slice(24, 0, hex\_instr)} already masks the
value down to bits \texttt{24:0} \emph{before} \texttt{Sign\_Extend} is called, so bit 25 of
that already-masked value is unconditionally 0 --- \texttt{pad\_ones} can never be true, and
negative (backward) call offsets were \emph{never} sign-extended.

\subsection*{Why this mattered}
Any function called from inside another function (not \texttt{main} itself), where the
callee is defined \emph{before} its caller in source order --- the completely ordinary C
pattern, e.g.\ a helper \texttt{f()} written above \texttt{main()} --- computes a
\emph{backward} jump. With the sign-extension bug, every such backward call computed
\texttt{target + 2\textsuperscript{25}} instead of \texttt{target}, landing in zero-filled
garbage memory. \texttt{f}'s body never executed at all; the caller's return register
simply held whatever value happened to already be sitting in it before the call --- a
silent, plausible-looking wrong answer, not a crash.

\textbf{Confirmed isolated to \texttt{\$call}, not systemic}: \texttt{DisassembleToBranchInstr}
(same file, line 302, listed above) does \texttt{Sign\_Extend(11, Get\_Slice(16,5,...))} on
its own 12-bit field (bits \texttt{16:5}, real sign bit at index 11) --- called with 11,
correctly indexed. This is consistent with \texttt{while}/\texttt{for} loops, which use
\texttt{\$branch} rather than \texttt{\$call} for their own backward jumps, having always
worked correctly throughout this project.

\subsection*{Verification}
\begin{itemize}
  \item Minimal repro: \texttt{int f(void) \{ return 6; \} int main() \{ return f(); \}}.
        Before the fix, \texttt{\$call f} resolved to \texttt{npc=0x2000018} and never
        entered \texttt{f}; after the fix, disassembly correctly shows \texttt{\$call l\_24},
        \texttt{f}'s body executes, and the program halts with \texttt{r1=6} exactly as
        expected.
  \item \texttt{test\_struct.c} (which calls helper functions defined above \texttt{main})
        went from wrong (\texttt{0xa}) to exactly correct (\texttt{0x0}) with this fix alone.
  \item Full regression: zero regressions among the other already-passing tests.
  \item \textbf{Currently in production}: \texttt{test\_struct} is part of the project's
        current continuously-passing baseline against the authoritative simulator binary,
        confirming this fix is live, not just a one-off verification build.
\end{itemize}

\section{Bug E3: instruction memory was half the documented size, silently truncating larger programs}

\textbf{Status}: \textbf{SOLVED} --- confirmed present in the current source.

\subsection*{The code, as currently fixed (\texttt{McodeClasses.hpp})}
\begin{lstlisting}[caption={\texttt{McodeClasses.hpp}, lines 139, 143 (current, fixed state)}]
uint32_t __instruction_memory[4096]; // 16KB (AparaReference.pdf p.6, Fig 1.1:
                                      // "16KB of instruction space", 4 bytes/instruction
                                      // -> 4096 words). Was [2048] (8KB) -- half the
                                      // documented size.
...
int Instr_Mem_Size_In_Words () {return (4*1024);}  // 16 KB, see note above
\end{lstlisting}

\subsection*{The bug (before this fix)}
\texttt{AparaReference.pdf}, p.~6, \S1, Figure~1.1, states explicitly: \emph{``The
instruction memory provides 16KB of instruction space to each accelerator. Each
instruction is 4-bytes.''} 16KB $\div$ 4 bytes = \textbf{4096 words}. The simulator's array
was sized \texttt{[2048]} --- half the documented capacity, carrying a comment that claimed
``16KB'' but provided only 8KB.

The instruction loader (\texttt{McodeAccelerator.cpp}'s \texttt{Init\_Instruction\_Memory})
silently \textbf{drops} --- logs an \texttt{Error:}, does not write, does not abort --- any
instruction whose program counter is \texttt{>= 2048}. A program larger than 2048 words runs
straight off the end of its own truncated body into zero-filled memory, with no crash and no
obvious symptom: execution simply continues into \texttt{\$null} until the simulator's tick
budget runs out, leaving the return register frozen on whatever value it last held, which
can coincidentally look like a plausible (but wrong) answer.

\subsection*{Why this mattered}
Two real test programs were large enough to trigger it: \texttt{test\_scalar\_full} (2688
words; 640 silently dropped) and \texttt{test\_spill} (2496 words; 448 dropped). Both were
initially misdiagnosed as having a remaining logic bug, when in fact their actual
return-value computation was simply never loaded into memory at all.

\subsection*{An intermediate false lead, corrected before finalizing}
A first attempt bumped the array further than strictly necessary (to \texttt{[16384]}) as a
verification-only change, and was briefly reverted on the strength of a verbal claim that
``2048 words is the real hardware limit.'' That claim was then checked directly against the
written ISA specification (cited above) and found to be incorrect --- 4096 words, not 2048,
is the documented figure. The fix shown above uses the spec-correct value, not the
over-sized verification value.

\subsection*{Verification}
\begin{table}[h]
\centering
\caption{Before/after, exact register values}
\begin{tabular}{@{}lccc@{}}
\toprule
Test & Before fix & After fix & Expected \\
\midrule
\texttt{test\_scalar\_full} & \texttt{0x7ff0} (stale leftover) & \texttt{0xc} & \texttt{0xc} (12) \\
\texttt{test\_spill} & \texttt{0x19} (stale leftover) & \texttt{0x1d1} & \texttt{0x1d1} (465) \\
\bottomrule
\end{tabular}
\end{table}
Combined with Bug~E2's fix, all three originally-broken tests (\texttt{test\_struct},
\texttt{test\_spill}, \texttt{test\_scalar\_full}) now produce exactly their expected
values. Full regression: zero regressions among the other 14 tests. \textbf{Currently in
production}: both tests are part of the project's current continuously-passing baseline.

% ============================================================================
\chapter{Unsolved Bugs}
\label{ch:unsolved}
% ============================================================================

These three bugs are confirmed, precisely root-caused, and reproduced --- but not fixed.
They are flagged here for the course staff/hardware team, since fixing C++ simulator
source is outside the scope of the Python C-compiler project these were found while
building.

\section{Bug E4 (major, correctness): \texttt{\$vreduce} on unsigned vector types sign-extends instead of zero-extending}

\textbf{Status}: \textbf{NOT SOLVED.}

\subsection*{The code (exact, current, unfixed)}
\begin{lstlisting}[caption={\texttt{McodeOperations.cpp}, lines 110, 146--198 --- \texttt{\_\_vreduce\_operation\_\_}}]
void __vreduce_operation__ (McodeOpcode sub_opcode, McodeType dest_type,
                             McodeType src_type, vector<uint64_t>& rs1_in_vector,
                             uint64_t& ovalue)
{
    // ... (lines 112-145: setup + the float-type branch, not relevant here)

    for(I=0, fI = rs1_in_vector.size(); I < fI; I++)
    {
        uint64_t ele = rs1_in_vector[I];

        // sign extend to 64-bits.
        int64_t r = (int64_t) Sign_Extend_64(src_type.Get_Nbits()-1, ele);  // <-- BUG: see below
        if(signed_flag)
        {
            int64_t r = (int64_t) Sign_Extend_64(src_type.Get_Nbits()-1, ele);
            if(sub_opcode == __ADD)
                s_result  = s_result + r;
            // ... (same pattern for MUL/OR/AND/XOR/XNOR/MAX/MIN)
        }
        else
        {
            if(sub_opcode == __ADD)
                result  = result + r;   // <-- uses the OUTER, sign-extended r
            // ... (same pattern for MUL/OR/AND/XOR/XNOR/MAX/MIN)
        }
    }
    ovalue = CastToU64 (signed_flag, dest_type.Get_Nbits(),
                         (signed_flag ? ((uint64_t) s_result) : result));
}
\end{lstlisting}

\subsection*{What's wrong}
The line marked first \texttt{// <-- BUG} unconditionally sign-extends \texttt{ele} into
\texttt{r}, \emph{before} the \texttt{if(signed\_flag)} branch is even reached --- so it runs
for every element regardless of the vector's actual declared signedness.
\begin{itemize}
  \item The \texttt{if(signed\_flag)} branch redeclares its own local \texttt{r}, identical
        to the outer one --- correct, just redundant.
  \item The \texttt{else} (unsigned) branch declares \textbf{no local \texttt{r} at all}.
        Every use of \texttt{r} inside it falls through to the outer,
        already-sign-extended \texttt{r}. The unsigned path never zero-extends \texttt{ele};
        it silently reuses the signed value computed for the wrong case.
\end{itemize}

\subsection*{Why this matters}
For an unsigned element with its top bit set (e.g.\ a \texttt{\$vu8} byte \texttt{0xFC} =
252), the correct behavior is to zero-extend it to \texttt{252} before summing. The bug
sign-extends it to \texttt{-4} instead --- the simulator computes the \emph{signed}
interpretation for a value the program declared \emph{unsigned}.

\subsection*{Confirmed reproduction}
\begin{lstlisting}[caption={Reproduction, exact numbers}]
vector (vu8, 8 elements): [1, 2, 3, 0xFC, 5, 6, 7, 8]
  $vreduce + on ($vi8)  -> 28    (correct, signed: 1+2+3+(-4)+5+6+7+8)
  $vreduce + on ($vu8)  -> 28    (simulator's actual output -- WRONG, should be 284)
                                  (correct, unsigned: 1+2+3+252+5+6+7+8 = 284)
\end{lstlisting}
Same pattern confirmed at \texttt{\$vu16} (expected \texttt{65538}, got \texttt{2}) and
\texttt{\$vu32} (expected \texttt{4294967295}, got \texttt{-1}) --- generic across all three
unsigned vector widths, not width-specific. Currently causes exactly 3 documented/expected
errors in this project's verification baseline, landing precisely on these three cases. A
standalone reproduction program independently confirms the same 3-error pattern against the
real simulator binary.

\subsection*{Suggested fix}
\begin{lstlisting}[caption={Minimal fix}]
else
{
    uint64_t r = ele & __mmask__(src_type.Get_Nbits());   // zero-extend, not sign-extend
    if(sub_opcode == __ADD)
        result = result + r;
    // ... (same for the remaining sub-opcodes)
}
\end{lstlisting}
\texttt{\_\_mmask\_\_} already exists and is used elsewhere in this same file, so no new
helper is needed.

\section{Bug E5 (major, correctness): 4-bit vector types (\texttt{vi4}/\texttt{vu4}) silently compute garbage}

\textbf{Status}: \textbf{NOT SOLVED} --- explicitly deprioritized, not used frequently in
this project's target workloads.

\subsection*{The code (exact, current, unfixed)}
\begin{lstlisting}[caption={\texttt{McodeOperations.cpp}, lines 50--74 --- \texttt{CastToU64}}]
uint64_t CastToU64 (int signed_flag, uint32_t nbits, uint64_t ival)
{
    uint64_t result;                 // <-- never initialized
    switch(nbits)
    {
        case 8:
            if(signed_flag) result =  ___signed_cast___(int8_t, ival);
            else  result =  ___unsigned_cast___ (uint8_t, ival);
            break;
        case 16:
            if(signed_flag)  result = ___signed_cast___(int16_t, ival);
            else  result =  ___unsigned_cast___ (uint16_t, ival);
            break;
        case 32:
            if(signed_flag) result =  ___signed_cast___(int32_t, ival);
            else  result =  ___unsigned_cast___ (uint32_t, ival);
            break;
        case 64:
            if(signed_flag)  result = ___signed_cast___(int64_t, ival);
            else   result = ___unsigned_cast___ (uint64_t, ival);
            break;
        // <-- no "case 4:", and no "default:" either
    }
    return(result);
}
\end{lstlisting}

\subsection*{What's wrong}
The \texttt{switch(nbits)} has cases for \textbf{8, 16, 32, 64 only}. For any 4-bit-wide
result --- every \texttt{vi4}/\texttt{vu4}/\texttt{vf4} vector element operation --- the
switch falls through with no case matching. \texttt{result} is declared but \textbf{never
assigned}, so the function returns whatever uninitialized garbage value happened to be
sitting on the stack at that point.

\subsection*{Confirmed reproduction (hand-computed, reproducible across repeated runs)}
\begin{table}[h]
\centering
\caption{4-bit vector add: expected vs.\ actual hardware result}
\begin{tabular}{@{}lccc@{}}
\toprule
Type & Expected (hand-computed) & Hardware actual & Result \\
\midrule
\texttt{vu8} (control) & \texttt{0x1112131415161718} & \texttt{0x1112131415161718} & exact match \\
\texttt{vi4} & \texttt{0x3333333333333333} & \texttt{0x0} & \textbf{WRONG} \\
\bottomrule
\end{tabular}
\end{table}
(\texttt{vi4} case: \texttt{a=0x1111111111111111}, \texttt{b=0x2222222222222222}, every
nibble \texttt{1+2=3}, no overflow possible --- yet the hardware returns exactly \texttt{0}.)

Type parsing itself is correct (the grammar correctly resolves \texttt{vi4\_t} to
\texttt{nbits=4}); the bug is purely the missing switch case, not type recognition.

\subsection*{Scope check --- audited, not guessed}
Every \emph{other} \texttt{switch(nbits)} block elsewhere in the simulator (in
\texttt{McodeNumeric.cpp} and five blocks in \texttt{McodeFpuUtils.cpp}) either explicitly
handles 4 bits, or fails loudly with \texttt{assert(0)} instead of silently returning
garbage. \texttt{CastToU64} is the \emph{only} one that declares its result variable
uninitialized with no \texttt{default:} case at all --- confirmed isolated to this one
function, not a systemic pattern across the codebase.

\subsection*{Suggested fix}
Add a \texttt{case 4:} doing a manual 4-bit sign-extend/mask (there is no native C++
\texttt{int4\_t} to reuse the existing \texttt{\_\_\_signed\_cast\_\_\_}/
\texttt{\_\_\_unsigned\_cast\_\_\_} macros with directly).

\subsection*{Status}
\texttt{vu4}/\texttt{vf4} are untested but share the identical code path, so should be
assumed equally broken until checked. This project's own instruction-coverage scope
explicitly excludes \texttt{vi4}/\texttt{vu4} for exactly this reason.

\section{Bug E6 (minor, diagnostic-only): misleading error messages in the post-condition file parser}

\textbf{Status}: \textbf{NOT SOLVED} --- does not affect execution correctness, only
diagnostic clarity.

\subsection*{The code (exact, current, unfixed)}
\begin{lstlisting}[caption={\texttt{McodeAccelerator.cpp}, \texttt{Verify\_Line}, relevant excerpts (lines 422--516)}]
else if (tokens[1] == "reg")
{
    //      reg <thread-id>  reg-id reg-value [reg-mask]
    if (tokens.size() < 3)
    {
        McodeRoot::Error ("Incomplete reg line in Verify_Line", this);   // line 427, correct here
        ret_val = 1;
    }
    ...
}
else if (tokens[1] == "mem")
{
    //      mem <address>    mem-value [mem-value-mask]
    if (tokens.size() < 4)
    {
        McodeRoot::Error ("Incomplete reg line in Verify_Line", this);   // line 473 -- says "reg", wrong branch
        ret_val = 1;
    }
    ...
}
else
{
    ret_val = 1;
    McodeRoot::Warning("Unknown verify line keyword " + tokens[0], NULL);  // line 515 -- checks tokens[1], reports tokens[0]
}
\end{lstlisting}

\subsection*{Bug E6a --- wrong keyword named in the ``incomplete mem line'' error (line 473)}
This branch is reached only when parsing a malformed \textbf{\texttt{mem}} line, but the
message says \textbf{``Incomplete reg line''} --- copy-pasted from the \texttt{reg}-handling
branch (line 427) a few lines above and never updated for its new home.

\subsection*{Bug E6b --- error message names the wrong token on an unrecognized line (line 515)}
Every branch in this function dispatches on \textbf{\texttt{tokens[1]}} (the file format is
\texttt{<thread-id> <keyword> <args...>}, e.g.\ \texttt{0 mem 0x80 0x1}). But the fallback
warning at line 515 prints \textbf{\texttt{tokens[0]}} (the thread-id) instead of the
keyword that actually failed to match. Feeding a line in a slightly different but
plausible-looking format (e.g.\ missing the leading thread-id, so what should be the
keyword shifts into \texttt{tokens[0]} and the real \texttt{tokens[1]} becomes something
else) lands in this branch and reports something like \texttt{Unknown verify line keyword
reg} --- which reads as though \texttt{"reg"} itself were an invalid keyword, when
\texttt{"reg"} \emph{is} valid; the real problem is its position in the line.

\subsection*{Suggested fix}
\begin{lstlisting}[caption={Two one-line fixes}]
McodeRoot::Error ("Incomplete mem line in Verify_Line", this);            // for E6a
McodeRoot::Warning("Unknown verify line keyword " + tokens[1], NULL);     // for E6b
\end{lstlisting}

\subsection*{Status}
Not fixed. Diagnostic-only --- does not affect any verification check's actual pass/fail
correctness, only the clarity of error messages encountered while reverse-engineering the
exact file format this parser expects.

% ============================================================================
\chapter{Related Findings (Not Source-Code Bugs)}
% ============================================================================

\section{Two disagreeing simulator binaries exist in the project tree}
A second, different \texttt{mcode\_run} binary exists elsewhere in the project tree,
confirmed (by MD5 checksum) to be a different build from the authoritative one used
throughout this project. It accepts a \emph{different} \texttt{PostCondition} file format
(no leading thread-id) and belongs to a separate RTL/testbench-oriented verification setup
--- not this project's compiler-targeted simulator. Despite having a \emph{newer} file
timestamp than the authoritative binary, file recency alone is not evidence of which build
is correct for this project; several pre-existing example result files elsewhere in the
tree appear to have been written against the non-authoritative binary's format and will
not verify correctly against the authoritative one.

\section{Two diverged simulator source checkouts caused transient confusion}
While testing Bugs~E2/E3's fixes, a local rebuild was discovered to diverge from the
historical baseline checkout in roughly 15 source files unrelated to either fix. This
caused two previously-passing tests to newly fail when run against the diverged
toolchain, even though they passed against the original checkout on the exact same
compiled object file --- not a regression from either fix, just evidence that the two
trees are genuinely different codebases, not ``the same engine plus two patches.''
Resolved by consolidating this project on a single authoritative binary; both affected
tests are part of the current passing baseline against that binary.

\section{Function pointers are blocked by a hard architectural limitation, not a bug}
There is no instruction that can load a function's absolute address into a register.
\texttt{\$call <label>} encodes a PC-relative offset, not an absolute address. \texttt{\$set}'s
immediate-field grammar (confirmed directly against the assembler's own parser source)
only accepts numeric literal tokens --- a label name in that position is a hard parse
failure, and the assembler has no label-resolution mechanism outside of control-transfer
instructions. This is categorically different from the bugs in this report: it is not
broken logic to fix, it is a missing capability that would need a new instruction or
assembler feature. The Python compiler's own IR/codegen support for function pointers is
fully built and ready; it is blocked purely by this assembler-level gap, and cannot be
worked around from the compiler side.

% ============================================================================
\chapter{Summary}
% ============================================================================

\begin{table}[h]
\centering
\caption{All six confirmed \texttt{engine\_isp} bugs, at a glance}
\small
\begin{tabular}{@{}cp{5.0cm}p{4.0cm}p{2.6cm}@{}}
\toprule
\# & Bug & File & Status \\
\midrule
E1 & \texttt{\$ld} sub-word loads ignored the real byte offset & DMEM load logic (upstream VM) & \textbf{SOLVED} \\
E2 & \texttt{\$call} disassembler sign-extended the wrong bit & \texttt{McodeDisassemble.cpp:266} & \textbf{SOLVED} \\
E3 & Instruction memory half the documented size & \texttt{McodeClasses.hpp:139,143} & \textbf{SOLVED} \\
E4 & \texttt{\$vreduce} unsigned sign-extends instead of zero-extends & \texttt{McodeOperations.cpp} (\texttt{\_\_vreduce\_operation\_\_}) & \textbf{NOT SOLVED} \\
E5 & \texttt{vi4}/\texttt{vu4} vector ops return garbage & \texttt{McodeOperations.cpp} (\texttt{CastToU64}) & \textbf{NOT SOLVED} \\
E6 & Misleading \texttt{Verify\_Line} error messages & \texttt{McodeAccelerator.cpp:473,515} & \textbf{NOT SOLVED} (diagnostic-only) \\
\bottomrule
\end{tabular}
\end{table}

Of six confirmed defects in the simulator's own source, three are fixed and verified
present in the currently-used authoritative binary (E1--E3), and three remain open,
confirmed, precisely root-caused, and reproduced, with suggested one-paragraph fixes
each (E4--E6). None of the three open bugs are fixable from the C-compiler side of this
project --- E4 and E5 require a C++ source change in the simulator itself, and E6 is a
pure diagnostics fix in the same codebase.

\end{document}
